% Generated draft: requires mathematical and source review.
\documentclass[11pt]{article}
\usepackage[margin=25mm]{geometry}
\usepackage{fontspec}
\setmainfont{Latin Modern Roman}
\usepackage{amsmath,amssymb,mathtools,unicode-math}
\setmathfont{Latin Modern Math}
\usepackage{enumitem,needspace,xurl,hyperref}
\setlength{\parindent}{0pt}
\setlength{\parskip}{5pt}
\setlength{\emergencystretch}{3em}
\newcommand{\N}{\mathbb{N}}
\newcommand{\Z}{\mathbb{Z}}
\newcommand{\Q}{\mathbb{Q}}
\newcommand{\R}{\mathbb{R}}
\newcommand{\C}{\mathbb{C}}
\newcommand{\F}{\mathbb{F}}
\newcommand{\A}{\mathbb{A}}
\newcommand{\PP}{\mathbb{P}}
\newcommand{\E}{\mathbb{E}}
\newcommand{\eps}{\varepsilon}
\newcommand{\dd}{\,\mathrm{d}}
\begin{document}
\section*{UnsolvedMath ID: 11000052}
\textbf{Source problem code:} catalogue entry
\subsection{Definitions and mathematical statement}
Let $\mathcal I_g$ be the Torelli group, the kernel of the homological action of the mapping class group of a closed orientable genus-$g$ surface. The extracted source writes a sequence of groups as $(\mathcal I_g)_k$, but does not supply its definition or indexing convention.

\textbf{Conjecture 5.5, as extracted.} For every integer $k\ge1$, the group denoted $(\mathcal I_g)_k$ is not finitely generated. Explicitly, for each such $k$ this means
\[
 \forall A\subset(\mathcal I_g)_k\text{ finite},\qquad
 \langle A\rangle\ne(\mathcal I_g)_k.
\]
The intended genus range and the definition of the subscripted series are essential missing data. The surrounding local source questions use $g\ge3$, but the conjecture itself does not repeat that range.

The research summary discusses the lower central series. For a group $G$, one common convention is $\gamma_1G=G$ and $\gamma_{k+1}G=[G,\gamma_kG]$. If $(\mathcal I_g)_k$ meant $\gamma_k\mathcal I_g$, the $k=1$ instance would assert that $\mathcal I_g$ itself is not finitely generated. This conflicts, for $g\ge3$, with the finite-generation statement explicitly recalled in the surrounding source, preserved in local ID 11000050. A shifted convention $G_0=G$, $G_{k+1}=[G,G_k]$ would instead start the stated range at the commutator subgroup. The excerpt does not establish that this shifted convention is intended; it is a possible resolution requiring source checking.

These groups must also be distinguished from the Johnson-filtration terms written $\mathcal I_g(k)$ in the adjacent records, including $\mathcal I_g(2)=\mathcal K_g$. Similar-looking notation is not enough to identify the two series. Nor do statements about their rational graded Lie algebras alone determine finite generation of the groups. The excerpt's final sentence introduces a following test question and does not add that missing question here.
\subsection{Short English statement}
The source conjectures non-finite-generation of every positive-index term of a Torelli-group series. Its series definition and indexing are missing; an unshifted lower-central reading would conflict with the surrounding finite-generation statement for the Torelli group itself.
\subsection{Sources}
\begin{itemize}
\item Local UnsolvedMath registry, permanent ID 11000052, source code catalogue entry; no individual proposer or proposal year supplied.
\item Farb, edited volume \emph{Problems on Mapping Class Groups and Related Topics}, Conjecture 5.5, PDF page 45; \url{https://www.math.uchicago.edu/~farb/papers/mcgbook.pdf}. The registry records PDF-font-structure extraction and access on 2026-07-29; the inventory distinguishes 329 labeled headings by source page.
\item Adjacent local IDs 11000050 and 11000051 record the $g\ge3$ finite-generation context and the distinct Johnson notation $\mathcal I_g(k)$ with $\mathcal K_g=\mathcal I_g(2)$. The source bibliography itself is not reproduced in this record.
\item The research summary mentions Hain and associated graded objects without precise primary references or a resolution of the indexing ambiguity.
\end{itemize}
\end{document}
